\chapter{Page Viewer 3D \& Analyse Volumétrique}
\label{chap:viewer}

La page \textbf{Viewer} (\texttt{viewer.html} et son script contrôleur \texttt{js/pages/viewer.js}) est le poste de commande scientifique et visuel de la plateforme. Conçue de manière granulaire, son interface reflète fidèlement les étapes d'une analyse d'image professionnelle, tout en masquant la complexité mathématique du rendu volumique WebGL sous-jacent.

\section{Le Moteur de Rendu Volumique (Fondation Mathématique)}

Le rendu volumique repose sur une approche de \textbf{Raymarching} direct programmée en GLSL. Le GPU calcule l'équation paramétrique du rayon optique $P(t) = \vec{O} + t \cdot \vec{D}$ et itère à travers une Bounding Box (AABB) encadrant le volume.

\subsection{L'Optimisation par Early Ray Termination}
Lors de l'accumulation front-to-back de l'opacité (\textit{Alpha Compositing}), l'algorithme tire parti de l'\textbf{Early Ray Termination} : si la densité lumineuse accumulée dépasse 95\%, le rayon s'arrête prématurément, ce qui réduit exponentiellement la charge graphique.

\subsection{Correction d'Opacité selon le Pas (Step Size Correction)}
Une des rigueurs mathématiques absolues implémentée dans notre algorithme concerne la conservation de l'énergie visuelle. Le nombre d'échantillons le long du rayon dépend de la taille du pas (\texttt{dt}). Si l'utilisateur ajuste la qualité de rendu (faisant varier \texttt{dt}), l'opacité $A$ lue dans la texture doit être rigoureusement corrigée pour que le nuage volumique ne paraisse pas soudainement plus dense ou plus transparent. La loi de Beer-Lambert est appliquée :
\begin{equation}
A_{\text{corrected}} = 1.0 - (1.0 - A)^{\frac{\texttt{dt}}{\texttt{dt\_reference}}}
\end{equation}
Cela garantit une exactitude scientifique visuelle, indépendamment de la puissance de la carte graphique de l'utilisateur.

\section{La Barre d'Outils Principale (Top Toolbar)}

La barre d'outils supérieure regroupe les interactions primordiales de la scène spatiale, divisées en trois catégories sémantiques : Tools, Layouts, et Visuals.

\subsection{Outils d'Interaction (Tools)}

\begin{itemize}
    \item \textbf{Navigate} : Outil par défaut basé sur le composant \texttt{OrbitControls}. Il permet la rotation (Orbital), le panoramique (Pan) et le zoom dynamique. Le pivot est automatiquement recalculé pour cibler le centre de masse du volume.
    \item \textbf{Slice (Coupe Dynamique)} : Active un plan de coupe manipulable par un Gizmo 3D (3 axes de rotation : Yaw, Pitch, Roll). Le shader GLSL reçoit l'équation de plan vectorielle et supprime (\texttt{discard}) dynamiquement les fragments via le produit scalaire des normales :
\begin{lstlisting}[language=C, caption=Clipping planaire en GLSL]
if (dot(fragPos, planeNormal) > planeDistance) { discard; }
\end{lstlisting}
    \item \textbf{Measure (Mesure 3D)} : Permet de placer des marqueurs topographiques. Grâce à l'algorithme de \textit{Depth-Picking}, le raycaster analyse le seuil d'opacité locale pour déterminer la surface du "nuage" de voxels. La distance euclidienne réelle est ensuite restituée grâce aux métadonnées physiques $V_s$ (en $\mu m$) :
\begin{equation}
d = \sqrt{\Delta x^2 + \Delta y^2 + \Delta z^2} \times V_s
\end{equation}
\end{itemize}

\subsection{Mises en Page et Vues (Layouts)}
\begin{itemize}
    \item \textbf{Presentation} : Bascule l'interface en plein écran et occulte tous les panneaux non essentiels.
    \item \textbf{Decompose by Channels} : Outil analytique majeur. Il duplique contextuellement le volume actif et répartit automatiquement chaque canal fluorescent dans une sous-fenêtre synchronisée. Ce mode permet d'isoler spatialement une co-localisation de protéines sans altérer la perspective.
    \item \textbf{Z-Stack Browser} : Active un mode d'inspection orthogonal "tranche par tranche" classique, désactivant la perspective 3D.
\end{itemize}

\subsection{Effets Visuels et Modificateurs (Visuals)}

\begin{itemize}
    \item \textbf{Render Mode (MIP vs DVR)} :
    \begin{itemize}
        \item \textit{MIP (Maximum Intensity Projection)} : Ce mode retient l'intensité lumineuse maximale absolue rencontrée par le rayon.
        \item \textit{DVR (Direct Volume Rendering)} : Mode photoréaliste qui simule l'accumulation physique de l'absorption lumineuse. Pour gérer la très haute dynamique (HDR) des pixels fluorescents accumulés, l'algorithme GLSL utilise la courbe de Tonemapping cinématographique \textbf{ACESFilm}, définie par l'approximation polynômiale suivante permettant d'éviter que le signal optique ne "brûle" à l'écran :
\begin{equation}
f(x) = \frac{x (a x + b)}{x (c x + d) + e}
\end{equation}
\textit{Avec $a=2.51$, $b=0.03$, $c=2.43$, $d=0.59$, et $e=0.14$.}
    \end{itemize}
    \item \textbf{Colormaps (LUTs)} : Offre des Tables de Correspondance alternatives (ex: Viridis, Inferno, Fire).
\end{itemize}

\section{Le Panneau de Contrôle Latéral (Left Menu)}

Le panneau de gauche offre un contrôle scientifique granulaire sur le traitement de l'image.

\subsection{Histogrammes et Gestion des Canaux (Channels)}
Pour chaque canal biologique présent, le panneau génère dynamiquement un histogramme statistique calculant la distribution de probabilité d'intensité $P(I)$.
L'interface de la carte de canal (Channel Card) est dotée d'une ergonomie très fine, regroupant une panoplie de contrôles agissant en temps réel sur la texture 1D de Fonction de Transfert (\textit{Transfer Function LUT}) du shader GLSL :

En se référant à l'interface, voici le détail des contrôles :
\begin{itemize}
    \item \textbf{(1) Checkbox Visibilité} : Permet d'afficher ou de masquer instantanément le canal sans décharger la texture VRAM.
    \item \textbf{(2) Sélecteur de Couleur} : Ouvre un Color Picker natif pour réassigner la teinte dominante du canal (qui module la LUT).
    \item \textbf{(3) Solo Channel (Icône Focus)} : Masque tous les autres canaux actifs pour se concentrer exclusivement sur celui-ci, indispensable lors du multiplexage lourd.
    \item \textbf{(4) Bouton Transparency (Icône Goutte)} : Modifie l'équation d'accumulation Alpha pour rendre ce canal spécifiquement plus diaphane.
    \item \textbf{(5) Curseurs de l'Histogramme (Sliding Handles)} :
    \begin{itemize}
        \item \textit{Curseur Gauche (Min)} : Tronque radicalement le bruit de fond optique (background noise).
        \item \textit{Curseur Droit (Max)} : Définit le plafond d'intensité lumineuse pour saturer les signaux forts.
        \item \textit{Curseur Central (Gamma $\gamma$)} : Applique une correction non-linéaire sur les tons moyens.
    \end{itemize}
    \item \textbf{(6) Boutons d'Auto-calibration} : (Auto, Soft, Contrast, Reset) appliquent des préréglages algorithmiques calculés sur la base de la distribution $P(I)$.
    \item \textbf{(7) Checkbox "Ignore Low"} : Si cochée, l'algorithme ignore le pic massif de pixels noirs lors du calcul de la courbe.
\end{itemize}

Dans le fragment shader, l'application de ces paramètres sur l'intensité brute $I$ d'un voxel s'exécute ainsi :
\begin{lstlisting}[language=C, caption=Lecture de la Transfer Function en GLSL]
// Normalisation de l'intensite brute entre [0.0, 1.0] selon les curseurs Min/Max
float normalizedI = smoothstep(thresholdMin, thresholdMax, intensity);
// Application du Gamma non-lineaire (Curseur Central)
normalizedI = pow(normalizedI, gamma);
// Lecture de la couleur (Sélecteur) dans la texture de la Colormap 1D
vec4 voxelColor = texture(lutMap, vec2(normalizedI, 0.5));
\end{lstlisting}

\subsection{Échelle Physique (Physical Scale)}

Le panneau d'échelle physique récupère le ratio \textit{Pixel to Micrometer} et dessine une barre de référence dynamique calibrée (ex: $50\,\mu m$) qui s'adapte géométriquement en temps réel lors du défilement de zoom de l'utilisateur.

\subsection{Qualité de Rendu et Optimisations Mémoire (Render Quality)}

C'est le système nerveux central garantissant la fluidité :
\begin{itemize}
    \item \textbf{Out-of-Core Bricking} : L'interface permet de définir le budget mémoire vidéo autorisé. Le volume charge progressivement ses tuiles 3D (Chunks) dans des \texttt{gl.TEXTURE\_3D}.
    \item \textbf{Dynamic LOD et Progressive Refinement} : Pour maintenir 60 FPS constants sur les embryons les plus lourds, le moteur écoute les événements de la caméra orbitale. Lors d'un mouvement, le ratio de pixels (Pixel Ratio) de l'écran est dynamiquement divisé (sous-échantillonnage) :
\begin{lstlisting}[language=Java, caption=Gestion du Dynamic LOD en JavaScript]
controls.addEventListener('change', () => {
    // Si la camera bouge, on degrade temporairement la resolution a 50%
    renderer.setPixelRatio(window.devicePixelRatio * 0.5);
    needsRefinement = true;
});

controls.addEventListener('end', () => {
    // A l'arret, restauration immediate de la nettete maximale
    renderer.setPixelRatio(window.devicePixelRatio);
    render();
});
\end{lstlisting}
\end{itemize}

\subsection{Environnement et Fond (Background)}

La flexibilité du canevas WebGL permet d'isoler visuellement les données (Noir, Blanc, et paramétrage du buffer WebGL avec \texttt{alpha: true} pour un export PNG sans fond transparent, vital pour les figures Illustrator).
